% sections/setup.tex -- Section 6 (paper_v2_vn, translated in L6 from paper_v2's sections/setup.tex;
% terms as in glossary_vn.tex). The Thiết lập paragraph gives the key settings (the supplementary
% document, supplement_vn.tex, has them all since N31), Table~\ref{tab:cost} the search's run
% times. Every number is \val{} from numbers.tex. Table: tables/t3_sets.tex (T3).
\section{Thiết lập thực nghiệm}
\label{sec:setup}

\paragraph{Các tập meta-path.} Chúng tôi so sánh các tập tìm được với các tập thủ công, cùng các
tập chọn ngẫu nhiên làm đối chứng, tổng cộng 24 tập (Bảng~\ref{tab:sets}).
\begin{itemize}
  \item \emph{Tìm được:} với mỗi kích thước tập, \code{a1}, \code{a2} và \code{c}
        (Mục~\ref{sec:fleet}).
  \item \emph{Thủ công:} M1--M8, thiết kế theo bốn nguyên tắc nêu rõ.
        (1) Mỗi loại bằng chứng một meta-path (mua kèm \code{PP}, đồng nhận xét \code{PUP}, chung
        danh mục \code{PCP}). (2) Ưu tiên meta-path ngắn. (3) Tránh meta-path gần rỗng (Mục~\ref{sec:data}), nên
        \code{PCCCP}, vốn đến được các danh mục anh em, được dùng thay cho \code{PCCP}. (4) Có các tập chính cùng kích thước
        với các tập tìm được, M5 và M6. Sáu tập còn lại, mỗi tập thử một ý tưởng, chẳng hạn chỉ hành vi (M3) hoặc chỉ
        danh mục (M4).
  \item \emph{Ngẫu nhiên (đối chứng):} năm tập cho mỗi kích thước, chọn đều và không lặp lại từ
        chính \val{search-cands} ứng viên của tìm kiếm, để biết việc chọn meta-path ảnh hưởng đến
        đâu.
\end{itemize}
Các tập thủ công và các tập chọn ngẫu nhiên được cố định trước khi có bất kỳ kết quả tìm kiếm hay
mô hình nào trên dữ liệu đã làm sạch. Các tập tìm được thì cố định trước mọi lần chạy HAN. Cả
\val{sets-distinct-paths} meta-path khác nhau trong 24 tập đều nằm trong \val{search-cands} ứng
viên, nên phương pháp tìm kiếm đã có thể chọn bất kỳ meta-path nào trong số đó.

% t3_sets.tex -- T3: the 24 path sets (paper_v2_vn, translated in L3 from paper_v2's
% tables/t3_sets.tex; HAN repository, arch.py, via numbers.tex). Since N31 it also lists the ten
% random draws, which were a table of their own in the appendix (t3b_random.tex, now in the
% supplementary document). The reasons of the manual sets
% follow LMSPS review/manual_metapath_sets.md (revised version).
\begin{widetable}[htbp]
  \centering
  \small
  \caption{Cả 24 tập meta-path: tìm được, thủ công và chọn ngẫu nhiên làm đối chứng. Một meta-path
    được viết từ sản phẩm đi ra
    và đại diện cho cả chính nó lẫn meta-path ngược của nó. Ví dụ, \code{PUP} nối hai sản phẩm được cùng
    một người dùng nhận xét. Tập có dấu * là tập chọn theo xác thực của HAN
    (Mục~\ref{sec:res-random}).}
  \label{tab:sets}
  \begin{tabularx}{\linewidth}{@{}llL@{}}
    \toprule
    Tập & Meta-path & Cách chọn \\
    \midrule
    \multicolumn{3}{@{}l}{\emph{Tìm được: LMSPS, \val{search-runs} lần chạy cho mỗi kích thước
      tập, kết hợp lại}} \\
    \code{a1}, 4 meta-path & \texttt{\val{set-f4a1}} & điểm trung bình tốt nhất, trong một cửa sổ
      \val{search-window-4} tập \\
    \code{a2}, 4 meta-path & \texttt{\val{set-f4a2}} & tập còn lại của cửa sổ đó \\
    \code{c}, 4 meta-path & \texttt{\val{set-f4c}} & tập đồng thuận của bỏ phiếu \\
    \code{a1}, 6 meta-path & \texttt{\val{set-f6a1}} & đơn giản nhất trong \val{search-window-6}
      tập của cửa sổ \\
    \code{a2}, 6 meta-path & \texttt{\val{set-f6a2}} & đơn giản thứ hai \\
    \code{c}, 6 meta-path & \texttt{\val{set-f6c}} & tập đồng thuận của bỏ phiếu \\
    \midrule
    \multicolumn{3}{@{}l}{\emph{Thủ công: thiết kế với sự hỗ trợ của AI theo các nguyên tắc nêu
      rõ, cố định từ trước}} \\
    M1 & \texttt{\val{set-m1}} & một meta-path hành vi, một meta-path nội dung \\
    M2 & \texttt{\val{set-m2}} & mỗi loại bằng chứng một meta-path: mua kèm, đồng nhận xét, danh
      mục \\
    M3 & \texttt{\val{set-m3}} & chỉ hành vi \\
    M4 & \texttt{\val{set-m4}} & chỉ danh mục \\
    M5 & \texttt{\val{set-m5}} & M2 và \code{PPUP}, dùng cho so sánh chính với 4 meta-path \\
    M6 & \texttt{\val{set-m6}} & M5 và hai meta-path rộng hơn, dùng cho so sánh chính với 6 meta-path \\
    M7 & \texttt{\val{set-m7}} & mọi meta-path dài tối đa hai bước, không chọn lọc \\
    M8 & \texttt{\val{set-m8}} & các meta-path dài qua người dùng, như các lần chạy tìm kiếm trước
      đó ưa thích \\
    \midrule
    \multicolumn{3}{@{}l}{\emph{Ngẫu nhiên (đối chứng): chọn đều và không lặp lại từ
      \val{search-cands} ứng viên, cố định từ trước}} \\
    \emph{lần chọn} & \emph{4 meta-path} & \emph{6 meta-path} \\
    1 & \texttt{\val{set-r4d1}} & \texttt{\val{set-r6d1}} \\
    2 & \texttt{\val{set-r4d2}} & \texttt{\val{set-r6d2}}\textsuperscript{*} \\
    3 & \texttt{\val{set-r4d3}} & \texttt{\val{set-r6d3}} \\
    4 & \texttt{\val{set-r4d4}} & \texttt{\val{set-r6d4}} \\
    5 & \texttt{\val{set-r4d5}} & \texttt{\val{set-r6d5}} \\
    \bottomrule
  \end{tabularx}
\end{widetable}


\paragraph{Độ đo.} Mỗi cạnh xác thực và thử nghiệm $(u, p)$ được xếp hạng so với
\val{split-negatives} mẫu âm, lấy mẫu một lần trong số các sản phẩm mà $u$ không có cạnh nối tới.     
Hạng $r$ của $p$ là số ứng viên có điểm cao hơn hoặc bằng điểm của $p$, tính cả chính $p$.
Như vậy, $p$ luôn được xếp sau các ứng viên bằng điểm với nó, nên các trường hợp bằng điểm
được tính bất lợi cho mô hình. Chúng tôi báo cáo Hit@10 ($[r \le 10]$) và NDCG@10
($[r \le 10] / \log_2(r + 1)$) trên \val{split-test} cạnh thử nghiệm. Chúng tôi không báo cáo
Recall@10 và Precision@10 vì với một sản phẩm liên quan cho mỗi lần xếp hạng, chúng lần lượt bằng
Hit@10 và Hit@10$/10$. Hit@10 lấy mẫu là độ đo chúng tôi đã cố
định từ trước, nhưng so với \val{split-negatives} sản phẩm ngẫu nhiên, một sản phẩm phổ biến dễ dàng
thắng: riêng mô hình phổ biến đã đạt \val{pop-hit}. Ngoài ra, các độ đo lấy mẫu có thể xếp thứ tự các mô
hình khác với xếp hạng toàn bộ \citep{krichene2022sampled}. Để kiểm tra, chúng tôi cũng xếp hạng
mỗi cạnh thử nghiệm so với mọi sản phẩm mà người dùng của nó không có cạnh nối tới (gần như toàn bộ
\val{after-products} sản phẩm), rồi báo cáo Hit@10 theo cạnh và Recall@20 theo người dùng, với mọi sản
phẩm thử nghiệm của một người dùng được xếp hạng cùng nhau, trên \val{test-users} người dùng có
cạnh thử nghiệm. Mô hình phổ biến được báo cáo bên cạnh các mô hình trong Bảng~\ref{tab:main}.
Các cạnh thử nghiệm cũng được chia theo số cạnh huấn luyện của người dùng (1--5, 6--20 và từ 21 trở lên).

\paragraph{Thiết lập.} Mỗi mô hình có một cấu hình duy nhất, cố định trước các lần chạy và dùng cho
mọi tập meta-path. Không có gì được tinh chỉnh theo từng tập hay theo kết quả thử nghiệm. Các thiết
lập được chọn để vừa với một GPU \val{env-gpu-mem}\,GB và để huấn luyện ổn định, chứ không được dò
tìm để tối đa độ chính xác. HAN có \val{han-set-num-heads} đầu chú ý \val{han-set-head-dim} chiều
cho mỗi meta-path và embedding người dùng \val{han-set-hidden} chiều. GAT có \val{gat-set-layers}
tầng \val{gat-set-hidden} chiều. metapath2vec học embedding \val{mtv-set-dim} chiều với cửa sổ
\val{mtv-set-window}. Mọi bộ tối ưu đều là Adam, với tốc độ học \val{mtv-set-lr} cho metapath2vec
và \val{han-set-lr} cho tìm kiếm cùng các mô hình còn lại. Việc huấn luyện dừng khi BPR xác thực
không cải thiện quá \val{han-set-min-delta} trong \val{han-set-patience} epoch
(với metapath2vec là \val{mtv-set-patience} lần kiểm tra). Sau đó, điểm lưu tốt nhất được thử
nghiệm một lần. Mọi lần chạy được báo cáo đều dừng sớm: HAN ở epoch
\val{han-stop-min}--\val{han-stop-max} trong tối đa \val{han-set-stage} epoch.

\paragraph{Seed và thống kê.} Phép chia được cố định, nên các seed chỉ thay đổi quá trình huấn
luyện. So sánh chính (\code{a1} so với M5 và M6, trong HAN và trong metapath2vec) và GAT chạy với
seed 1--3, các tập khác với seed~1. Qua ba seed, Hit@10 thử nghiệm của các lần chạy chính dao động
nhiều nhất \val{main-hit-sd-max} (độ lệch chuẩn), nên các chênh lệch nhỏ giữa những tập chạy với
một seed nằm trong nhiễu giữa các seed. Vì vậy, chúng tôi không xếp hạng chi li những tập như vậy. Chúng
tôi so sánh các tập chính bằng kiểm định $t$ hai phía của Welch trên ba seed của mỗi bên
\citep{welch1947generalization}. Với chỉ ba seed mỗi bên, kiểm định này khó phát hiện các chênh lệch nhỏ, nên
chúng tôi cũng nêu mỗi chênh lệch theo đơn vị độ lệch chuẩn giữa các seed. Chỉ Hit@10 lấy mẫu được cố định từ
trước. Các độ đo khác và giá trị $p$ của chúng chỉ mang tính bổ trợ. Các giá trị $p$ này không được hiệu
chỉnh cho việc thực hiện nhiều phép so sánh. Quan hệ giữa các tập dùng hệ số tương quan Spearman với giá
trị $p$ hoán vị từ \val{perm-shuffles} lần xáo trộn. \emph{Tập chọn theo xác thực} là tập có BPR
xác thực thấp nhất.

\paragraph{Phần cứng và chi phí.} Mọi lần chạy đều dùng một GPU NVIDIA \val{env-gpu} với Python
\val{env-python} và PyTorch \val{env-torch} (CUDA \val{env-cuda}). Toàn bộ \val{search-runs-total}
lần chạy tìm kiếm mất khoảng \val{search-total-h}~giờ~GPU (Bảng~\ref{tab:cost}), còn các lần chạy mô
hình mất khoảng \val{models-hours}~giờ~GPU.
